% Transcription — fonds Alexandre Grothendieck, shelfmark 69, batch 1 (pages 1-20).
% Established from the facsimile archives/batches/69/batch-01.pdf, cut from a
% copy of 69.pdf fetched through the project's relay
% (https://grothendieck-relay.onrender.com/source/69.pdf), per the skill
% transcribe-grothendieck. The folder is 119 pages in six batches (1-20,
% 21-40, 41-60, 61-80, 81-100, 101-119), transcribed in parallel; this is the
% first. The raw scan's cover sheet was removed from the batch PDFs, so PDF
% page k is archivists' page k; the pencilled numbers bottom-left (« 14 » ...
% « 20 » checked on the leaves) agree.
%
% Pass: Opus 5.5 (claude-opus-5-5), 2026-09-27 — first pass, unchecked against the pages by a human.
%
% Pages skipped: 1, a tan cover sheet bearing only a pencilled « (34p) » top
% right (a page count, not a heading, and not certainly his), so it gets no
% \page marker and no heading. Pages summarised: none. Pages 2-20 are his
% working notes in blue ink on plain white paper, all transcribed.
%
% His pagination: the leaves are numbered by him, top centre, mostly circled:
% 1-10 (pages 2-11), « 10 bis » (12), 11-13 (13-15), « 13 bis » (16), 14-17
% (17-20). Page 10 ends « TSVP ». Each page opens with a \note giving his
% number. The run presumably continues into batch 2 (page 20 ends on a
% finished proof and a NB, so the break is clean).
%
% What the batch is about. The 27-vertex graph of the « graphe cubique »,
% i.e. the configuration of the 27 lines on a cubic surface, characterised
% combinatorially, and its automorphism group W(E_6).
%   2-10   Th.: a graph S with card S = 27, every vertex of valency 10, and
%          every edge in exactly one triangle is (and is the only) « graphe
%          cubique ». Proof through Lemmas 1-9: a triangle T = {u, v, w},
%          the partition S* = E(u)* ⊔ E(v)* ⊔ E(w)* (8 each), the involution
%          σ_T (an automorphism), the quotients P_u = E(u)*/σ with the maps
%          ω_v, δ_v into F_2(P_v), the bijections c_{wu} on the 2.2-partitions
%          Π_22(P_u) and their transitivity (Lemma 6), the set Γ of the 16
%          triangles of S* modulo σ as a torsor under Ψ(V) = {x+y+z = 0}
%          (Lemma 7), then an « Épilogue » labelling the 24 points α_i, β_i,
%          γ_i, α'_i ... and the incidence tables (I), (I'), (II), (II'),
%          (III), (III'), ending « S est déterminé ... à isom. près ».
%   11-12  Counting: the chain of subgraphs D_0 ⊂ ... ⊂ D_7 (point, edge,
%          triangle, « balise », « préprisme », prism, « bitriangle »), their
%          embeddings Pl_i as homogeneous spaces under W = Aut C, cardinals
%          27, 270, ..., 8640, configurations 1, 27, 135, 45, 1080, 4320,
%          720, 120, and card W = 2^7.3^4.5 = 51840.
%   13     An explicit model: S = I × {0,1,2} ⊔ I × Π × ω(I) for two
%          3-element sets I, Π (ω(I) the circular orders), with the rules of
%          linkage, and the action of σ.
%   14-15  Proposition on a central extension E_0 of an abelian group G by
%          an F_2-vector space ε, the Baer product over Ψ(G) = Ker(G^3 -> G),
%          the lifts ρ_i, the quadratic map q : V -> ε; conditions (i)-(iv)
%          equivalent when 2G = 0, and existence iff rg V <= 2.
%   16     « 13 bis »: the general axioms (a) unique third vertex, (b) regular
%          of valency ν; Lemma 1: ν even, card S >= 3(ν-1), equality iff every
%          vertex outside a triangle is linked to exactly one of its vertices;
%          side computations with ν = 2(μ+1); small sketches.
%   17-20  « Graphes cubiques C : sommet marqué u »: C with a marked vertex
%          is the data of a 5-set Q_u (the triangles through u) and a torsor
%          P under Ψ = Ker(F_2^{Q_u} -> F_2); rules a)-d) rebuilding the
%          graph; a Lemma (x, y linked iff exactly one triangle t through u
%          with t - E(x) = t - E(y)); Corollary: W_u is an extension of S_5
%          by F_2^4, i.e. the Weyl group W(D_5), of index 27 in W ≃ W(E_6);
%          NB on « bases » of C (40 non-linked pairs in E(u), a unique base b
%          through a given pair), and the positions of u relative to a base
%          b and its opposite b': 16 + 16 + 40 = 72.
%
% Notation, fixed once. E(s) is his set of neighbours (« liés »), E(u)* =
% E(u) - {v, w}, S* = S - T. His P-shaped letter for the 3-element set of
% page 7 is set \sqsupset, with a note. His script P for 2-subsets (p. 5) is
% \mathfrak{P}. The set of triangles on p. 7-8 is written with a letter read
% as Γ (\tilde{\Gamma}, \Gamma), which he reuses on p. 14 for Ψ(G). On
% pages 17-19 the 5-set is Q_u (the page writes it in a script that is
% sometimes nearer a P; noted once). Double-barred 1 is \mathbf{1}. The
% symmetric group is \mathfrak{S}. His « de = » (de même) is kept as written.
%
% Hardest pages: 3 (dense insertions and a cancelled Lemma 4 with a drawing),
% 9-10 (the incidence tables, whose primes are uncertain), 18 (the Lemma's
% interleaved rewritings) and 20 (a large cancelled block, mostly \ill{}).
% Readings to check first: the primes in tables (I'), (III) and (III');
% « balise » and « préprisme » on page 11; « Compte » at the foot of page 12;
% « produit de Baer » on page 14; the margin note on page 20.
\documentclass[11pt,a4paper]{article}
\input{../preamble/grothendieck}

\folder{69}
\batch{1}
\pages{1}{20}
\dating{[à partir de 1976]}
\watermark{Édition de démonstration}
\foldertitle{Graphes cubiques : notes manuscrites (s.d.), lettre (s.d.).}

\begin{document}

\page{2}
\note{il numérote lui-même ses feuillets en haut, au centre : celui-ci porte
« 1 », et la suite court jusqu'à « 17 » (page 20), avec deux feuillets
« bis » (pages 12 et 16)}

\emph{Th.} Soit $S$ un graphe tel que

a) $\operatorname{card} S = 27$

b) $\forall s \in S$, $\operatorname{card} E(s) = 10$ \quad ($E(s) = \lbrace t
\in S \mid t \text{ lié à } s \rbrace$)

c) $\forall s, t \in S$, $s, t$ liés, $\exists!\, u \in S$, $u$ lié à $s$ et $t$.

Alors (et alors seulement) $S$ est un \emph{graphe cubique}.

\emph{Dém.} Soient $u \in S$\struck{, $v \in E(s)$} ($\operatorname{card} S
\geqslant 27$ !) $v \in E(u)$ i.e.\ $v$ lié à $u$ ($\operatorname{card} E(u) =
10$), $w \in S$ tel que $w$ lié à $u$ et $v$ (axiome c)). \add{On dit que $T =
\lbrace u, v, w \rbrace$ est un « triangle ».}

\emph{Lemme 1.} Tout $s \in S - \lbrace u, v, w \rbrace$ \add{$= S^{*}$} est lié
à un et un seul des él.\ $u, v, w$ [i.e.\ $S^{*} = S - \lbrace u, v, w \rbrace$
est \emph{réunion disj.} de $E(u)^{*} = S^{*} \cap E(u)$, $E(v)^{*}$, $E(w)^{*}$
(chacun de cardinal 8)].

Comme $\operatorname{card} E(u) = 10$, \struck{tiens} \ldots\ $\operatorname{card}
E(u) \cap S^{*} = 8$ car $E(u) \cap S^{*} = E(u) - \lbrace v, w \rbrace$, de même
$\operatorname{card} E(v) \cap S^{*}$, $E(w) \cap S^{*}$ \add{($E(v)^{*}$,
$E(w)^{*}$)} sont de cardinal 8. Or (par c)), $E(u) \cap E(v) = \lbrace w
\rbrace$ donc $E(u)^{*} \cap E(v)^{*} = \emptyset$ et de même pour les autres 2
couples, donc $E(u)^{*}, E(v)^{*}, E(w)^{*}$ disjoints, donc la
\struck{\ill{}} \add{réunion} est de cardinal $3 \times 8 = 24$. Comme $S^{*}$
est de cardinal $27 - 3 = 24$ par a), on trouve que $S^{*}$ est réunion
(disjointe) des $E(u)^{*}, E(v)^{*}, E(w)^{*}$, cqfd.

\emph{Lemme 2.} $\forall x \in E(u)^{*}$, $\exists!$ \uncertain{$x'$} $\in
E(v)^{*}$ lié à $x$. (De même pour $v, w$).\note{le $E(u)^{*}$ de l'énoncé est
récrit sur un premier $E(\ldots)$ chargé d'encre}
En effet, \struck{l'unicité \add{des liés}} reformulation de c).

\emph{Corollaire.} $\exists!\, \sigma = \sigma_T : S \to S$ tel que $\sigma_T$
\uncertain{laisse inv.} $u, v, w$ et que $\forall x \in S^{*} = S - \lbrace u,
v, w \rbrace$, $\sigma(x)$ soit lié à $x$ et $x$ et $\sigma_T(x)$ liés au
\uncertain{même} élément $u, v, w$ de $T$.\note{« même » est écrit comme un
signe $=$}
\drawing{en marge gauche, un triangle de sommets $w$ (en haut), $u$ et $v$ (en
bas), et un point $x$ relié à $u$ par un trait}

Ainsi, dès que $E(u)^{*}, E(v)^{*}, E(w)^{*}$ se décomposent \uncertain{par} 4
couples de \add{pts liés}.

\emph{Lemme 3.} Soit $x \in E(u)^{*}$. Alors $\operatorname{card} E(x) \cap
E(w)^{*} = 4$.

Si $x \in E(u)^{*}$, $y \in E(w)^{*}$, et si $x$ et $y$ liés, $x$ et
$\sigma_T(y)$ non liés (et inversement.) En d'autres termes, $E(x) \cap
E(w)^{*}$ contient un élément et un seul de chaque paire de sommets liés de
$E(w)^{*}$.

\emph{Dém.} Si $x$ était lié à $y$, il ne peut être lié à $\sigma_T y$, sinon
$\lbrace y$ et $\sigma_T(y) \rbrace$ seraient \uncertain{côté} de deux
triangles distincts $\lbrace w, y, \sigma_T(y) \rbrace$ et $\lbrace x, y,
\sigma_T(y) \rbrace$. \struck{Il faut \ill{} \add{de $x$ et $\sigma_T(y)$}
\ill{} les \ill{}}.

\page{3}
\note{« 2 » en haut au centre}

Considérons
\[
A_x = E(x) \cap \lbrace S^{*} - \sigma_T(x) \rbrace = E(x) \cap (E(v)^{*} \cup
E(w)^{*}) = E(x) - \lbrace u, \sigma(x) \rbrace
\]
son cardinal est $10 - 2 = 8$ par b), \struck{supposons que} : il se décompose
en réunion disjointe de $E(x) \cap E(v)^{*} = A_v$ et $E(x) \cap E(w)^{*} =
A_w$, supposons que l'un \struck{des deux \ill{}} \add{$A_v$, $A_w$ (disons
$A_w$)} de cardinal $\geqslant 4$, p.ex.\ $> 4$ (sinon on intervertit $v$ et
$w$). Alors $\exists\, y, y'$ liés dans $E(w)^{*}$ qui sont liés \ill{} $x$,
absurde \ill{}

\emph{Corollaire.} Si $x \in E(u)$, $y \in E(w)$ sont liés, $\sigma_T(x)$ et
$\sigma(y)$ sont liés (et inversement.)

Si $x$ lié à $y$, $x$ non lié à $y' = \sigma_T(y)$, donc $\sigma_T(x)$ lié à
$y'$, cqfd.

\emph{Corollaire.} $\sigma_T = \sigma$ est un automorphisme de $S$.

\marginal{NB}
\struck{\emph{Lemme 4.} Soient $x, y \in E(u)^{*}$ non liés. Alors
$\operatorname{card} E(x) \cap E(y) \cap E(w)^{*} = 2$. i.e.\
$\operatorname{card} \omega_x$ \ldots\ \emph{Dém.} S'il était $\neq 2$, il
\ill{} $> 2$ \ill{} par $\sigma(y)$.}\note{tout ce passage est barré en
travers et recouvert de boucles ; un dessin, barré lui aussi, s'y mêle}
\drawing{barré avec le Lemme 4 : un triangle de sommets $w$, $u$, $v$, dont
les côtés se prolongent en trois demi-droites marquées $\gamma_4$, $\beta_4$,
$\alpha_4$}

\add{Introduisons $P_u = E(u)^{*}/\sigma$ \ill{} et $\omega_v(x) =$ image de
$E(x) \cap E(v)^{*}$ dans $P_v$, donc $\omega_v(\sigma x) = \complement
\omega_v(x)$.}\note{ajout serré dans la marge de droite ; le signe devant
$\omega_v(x)$ est lu comme le complémentaire}

Soient $x, y \in E(u)^{*}$, considérons
\[
\omega_v(x) = E(x) \cap E(v)^{*}, \qquad \omega_v(y) = E(y) \cap E(v)^{*}
\]
\struck{et comme} donc on peut considérer $\omega_v(y)$ et $\omega_v(x)$ comme 2
sections de $E(v)^{*}$ sur $P_v = E(v)^{*}/\mathbb{F}_2$, dont leur différence
est un élément
\[
\delta_v(x, y) \in \mathbb{F}_2(P_v)
\]

\emph{NB} $\delta_v(x, \sigma(y)) = \delta_v(\sigma x, y) = \delta_v(x, y) +
\mathbf{1}_v$\note{un 1 à double barre} (\uncertain{fonction} cte égale à 1 sur $P_v$)

Donc
\[
\left\lbrace
\begin{aligned}
\operatorname{supp} \delta_v(x, y) &= \lbrace \beta \in P_v \mid \omega_v(x) \cap
\beta \neq \omega_v(y) \cap \beta \rbrace \\
P_v - \operatorname{supp} \delta_v(x, y) &= \lbrace \beta \in P_v \mid
\omega_v(x) \cap \beta = \omega_v(y) \cap \beta \rbrace
\end{aligned}
\right.
\]
$=$ image de $\omega_v(x) \cap \omega_v(y)$ dans $P_v$.

\emph{Lemme 4.} \struck{$\operatorname{Card} \operatorname{supp} \delta_v(x, y)
= 2$,} Supposons $x, y \in E(u)^{*}$, $x$ et $y$ \emph{non liés}, alors
$\operatorname{card} \operatorname{supp} \delta_v(x, y) = 2$ [i.e.\
$\operatorname{card} \omega_v(x) \cap \omega_v(y) = 2$, où $\omega_v(x) \cap
\omega_v(y) = E(x) \cap E(y) \cap E(v)^{*}$].

Sinon, ce cardinal serait $> 2$ ou $< 2$, et quitte à remplacer $y$ par
$\sigma(y)$, on peut supposer \ill{}

\page{4}
\note{« 3 », cerclé, en haut au centre}

qu'il est $> 2$. Soient \struck{donc} \add{d'abord} $s, t \in E(w)^{*}$ liés
\add{tous deux} \ill{} lié à $v$.

\emph{Lemme 5.} a) Soient $x \in E(u)^{*}$, $s \in E(w)^{*}$, $x$ et $s$ liés.
Soit $\xi \in$ \struck{$E(\ldots)$} lié à $x$ et $s$ \add{($\xi$ lié \ill{})}.
Alors $\xi \in E(v)^{*}$.

b) Soient $x, y \in E(u)^{*}$, $s, t \in E(w)^{*}$, $x$ et $y$ liés à $s$ et à
$t$. [donc par lemme 3, $x$ et $y$ non liés, $s$ et $t$ non liés]. \struck{Soit}
Soient $\xi, \xi'$ tels que \struck{$\xi$ $(x, s)$ et $(y, t)$} $\xi$ lié à $x$
et $s$, $\xi'$ lié à $y$ et $t$. Alors $\xi, \xi' \in E(v)^{*}$, et $\xi, \xi'$
liés.

\emph{Démonstration du lemme 5.} \struck{Soient} a) On a $\xi \in S^{*}$ car
$\xi \neq u, v, w$ (car $u$ non lié à $s$, $v$ non lié à $x$, $w$ lié ni à $x$
ni à $s$) et $\xi \notin E(u)^{*}$, sinon on aurait $\xi = \sigma(x)$ et $\xi$
non lié à $s$ par lemme 3, de $=$\note{« de même »} $\xi \notin E(w)^{*}$, donc
$\xi \in E(v)^{*}$.

b) Comme $\lbrace x, \xi, s \rbrace$ est un triangle, $\xi'$ est lié à un
\add{sommet} \struck{et} (et un seul) de ce triangle. Or il n'est pas lié à $x$
(sinon $\lbrace \xi', t \rbrace$ seraient côtés de deux triangles distincts
$\lbrace \xi', t, x \rbrace$ et $\lbrace \xi', t, y \rbrace$) ni à $s$
(\struck{\ill{}} raison symétrique) donc il est lié à $\xi$.
\drawing{à gauche du texte, un quadrilatère $x$, $t$, $y$, $s$ (côtés $xt$,
$ty$, $ys$, $sx$) ; de $x$ descend un triangle pointant vers $\xi$, de $t$ un
autre pointant vers $\xi'$}

\emph{Dém.\ du lemme 4.} Supposons $x, y \in E(u)^{*}$ \add{($x \neq y$)} liés à
$s, t, r \in E(w)^{*}$ \add{distincts}.
Considérons le triangle $\lbrace x, s, \xi \rbrace$, donc (lemme 5 a)) $\xi \in
E(v)^{*}$, et si $\xi' = \sigma(t) \in E(v)^{*}$ \add{on applique à $(x, y)$,
$(s, t)$} (d'après le lemme 5 b)) $\lbrace$\uncertain{$y$}$, t, \xi' \rbrace$ est
un triangle, mais aussi $\lbrace y, r, \xi' \rbrace$ en appliquant à $(x, y)$
$(s, r)$. Donc $\lbrace y, \xi' \rbrace$ \uncertain{appartient} à 2 triangles,
absurde.\marginal{dans la marge, à hauteur de ces lignes : « $x$ \quad $y$
\qquad $s$ \quad $t$ \quad $r$ »}

\struck{Lemme 6} \add{Corollaire.} Il existe une application \add{bijective,
canon.} et une seule de $\Pi_{22}(P_u)$ dans $\Pi_{22}(P_w)$ ($\Pi_{22}$
indique partitions de type $2.2$) telle que pour \add{tout} couple $x, y$ de
pts \struck{\ill{} à 2 \ill{}} \struck{\ill{}} liés de $E(u)^{*}$,
\struck{la partition \ill{}} on ait
\[
c_{w,u}\bigl(\lbrace \pi_u\lbrace x, y \rbrace, \complement_{P_u} \pi_u\lbrace
x, y \rbrace \rbrace\bigr) = \lbrace \operatorname{supp} \delta_v(x, y),
\operatorname{supp}(\delta_v(x, y) + \mathbf{1}_{P_v}) \rbrace .
\]
\marginal{$\pi_u : E_u^{*} \to P_u = E_u^{*}/\mathbb{F}_2$}
\note{le mot qui précède « liés » est pris dans les lignes biffées ; le
lemme 4 porte sur des points \emph{non} liés}

\page{5}
\note{« 4 », cerclé, en haut au centre}

Il faut montrer que le 2\textsuperscript{ième} mb ne dépend que de la
partition $\in \Pi_{22}(P_u)$, pas du choix de $(x, y)$ ! Pour ceci, on note

a) Elle est indépendante par permutation de $x$ et $y$ \struck{\ill{}} (car
$\delta_v(x, y) = \delta_v(y, x)$)

b) Elle ne change pas si on change $x$ et $\sigma(x)$ (donc aussi $y$ en
$\sigma(y)$, ou $x$ en $\sigma(x)$ et $y$ en $\sigma(y)$) (car $\delta_v(\sigma
x, y) = \delta_v(x, y) + 1$)

c) Si $x, y$ sont fixés, \struck{posons} \add{prenons} \struck{$s \neq t$}
$\lbrace s, t \rbrace =$ \struck{$E(x) \cap E(y) \cap E(\ldots)$} \add{(lemme
4)} $E(x) \cap E(y) \cap E(w)^{*}$, donc \ill{} $\lbrace x, y \rbrace = E(s)
\cap E(t) \cap E(u)^{*}$, \struck{soient} \add{liés} $\lbrace x', y' \rbrace =
E(s) \cap E(\sigma(t)) \cap E(u)^{*}$, alors $\pi_u(\lbrace x', y' \rbrace) =
\complement_{P_u} \pi_u(\lbrace x, y \rbrace)$, \struck{donc $x', y'$} et
$\lbrace x', y' \rbrace$ donne donc lieu à la même partition $\in
\Pi_{2,2}(P_u)$ que $\lbrace x, y \rbrace$, mais à $\lbrace x', y' \rbrace$ la
recette \struck{prise} de l'énoncé du corollaire associe la partition en
$\pi_v(E(x') \cap E(y') \cap E(v)^{*})$ [$= \lbrace s, \sigma(t) \rbrace$] et
son complémentaire, qui \struck{\ill{}} est égale à $\pi_v\lbrace s, t \rbrace$
[$= \operatorname{supp}(\delta_v(x, y) + 1)$] et son complémentaire \ldots

Énoncé plus fort :

\emph{Corollaire 2.} Soit $\mathfrak{P}'_2(E_u^{*})$ l'ens.\ des
\struck{couples} \add{parties à} 2 él.\ distincts de $E_u^{*}$ formées de deux
él.\ non \emph{liés}.\note{la lettre qu'il écrit ici est une grande P
cursive, rendue $\mathfrak{P}$ dans tout le lot}
Soit $R_{u,v}$ la relation \struck{\ill{}} dans $\mathfrak{P}'_2(E_u^{*})
\times \mathfrak{P}'_2(E_v^{*})$ :
\[
R_{u,v}(\alpha, \beta) \overset{\text{déf}}{\iff} \text{tt él.\ de } \alpha
\text{ est lié à tt él.\ de } \beta .
\]
Alors $R_{u,v}$ \add{\uncertain{définit}} une bijection de
$\mathfrak{P}'_2(E_u^{*})$ sur $\mathfrak{P}'_2(E_v^{*})$, soit
\struck{$\varphi$}
\[
\tilde{\varphi}_{vu} : \mathfrak{P}'_2(E_u^{*}) \to \mathfrak{P}'_2(E_v^{*})
\]
\[
\tilde{\varphi}_{vu}(\sigma \tilde{\alpha}) = \sigma
\tilde{\varphi}_{vu}(\tilde{\alpha}), \qquad \pi_v \tilde{\varphi}_{vu}(\lbrace x,
\sigma y \rbrace) = \complement_{P_v} \pi_v \tilde{\varphi}_{vu}(\lbrace x, y
\rbrace)
\]
\marginal{en oblique dans le coin inférieur gauche : \struck{Énoncé} plus
simple avec $R'_{u,v}(\alpha, \beta)$ : tt él.\ de $\alpha$ \emph{non lié} à
tt él.\ de $\beta$, [$\iff R_{u,v}(\sigma\alpha, \beta) \iff R_{u,v}(\alpha,
\sigma\beta)$] et car les applications $\varphi'_{vu} :
\mathfrak{P}_2(E_u^{*}) \to \mathfrak{P}_2(E_v^{*})$ forment un syst.\
\emph{\uncertain{transitif}} d'iso (compris \ill{} lemme 6 ci-dessous)}

\page{6}
\note{« 5 », cerclé, en haut au centre}

donc la partition définie par \struck{$\pi_u(x, \sigma y$} $\pi_v
\tilde{\varphi}_{vu}(\lbrace x, y \rbrace)$ ne change pas en remplaçant $x$ par
$\sigma(x)$, ou $y$ par $\sigma(y)$, ou les deux à la fois) donc ne dépend que
de \struck{$\pi_u \varphi$} $\alpha = \pi_u(\lbrace x, y \rbrace) \in
\Pi_{2,2}(P_u)$ ; enfin, elle ne change pas si on remplace $\alpha$ par
$\complement_{P_u} \alpha$ \ldots
\drawing{en haut à gauche, une rangée de petits schémas : des segments aux
extrémités marquées $x\ x'$, $y\ y'$, $\xi\ \xi'$, $\eta\ \eta'$, $s\ s'$,
$t\ t'$, surmontés d'arcs en chevrons qui relient les segments, et sous
eux des colonnes $\binom{s}{t}$, $\binom{s'}{t'}$, $\binom{s}{t'}$,
$\binom{s'}{t}$, $\binom{x}{y}$, $\binom{x'}{y'}$, $\binom{x}{y'}$,
$\binom{x'}{y}$ ; un trait vertical sépare les deux premiers groupes, et le
dernier groupe (avec $u'$, $v'$) est couvert de hachures}

\struck{Corollaire} \add{Lemme} 6. $c_{wv} c_{vu} = c_{wu}$ --- donc les
$c_{..}$ définissent un syst.\ transitif d'iso entre les $\Pi_{22}(P_u)$,
$\Pi_{22}(P_v)$, $\Pi_{22}(P_w)$.

\struck{Soit $\lbrace s, t \rbrace = c_{vu}\lbrace x, y \rbrace$ \add{($x, y
\in E(u)^{*}$ non liés)}, soit $\xi \in E(w)^{*}$ lié à $x$ et $s$, (lemme 5
a)), $\eta \in E(w)^{*}$ lié à $y, t$.}
\drawing{en marge, les points $w$ (en haut), $u$, $v$, et un angle droit
de sommet $s$ ; $x$ et $y$ à gauche, $t$ à côté de $s$}

Cela \struck{résulte} \add{est inclus dans le} résultat plus précis

\emph{Corollaire.}
\[
\tilde{\varphi}_{wv} \tilde{\varphi}_{vu}(\tilde{\alpha}) = \sigma\bigl(
\tilde{\varphi}_{wu}(\tilde{\alpha}) \bigr) \quad \text{i.e.} \quad
\tilde{\varphi}_{wv} \tilde{\varphi}_{vu} = \sigma \tilde{\varphi}_{wu}
\]
i.e.
\[
\left.
\begin{array}{ll}
x, y \in E(u)^{*} & x, y \text{ liés à } s, t \\
s, t \in E(v)^{*} & s, t \text{ liés à } \xi, \eta \\
\lbrace \xi, \eta \rbrace \in E(w)^{*} &
\end{array}
\right\rbrace \Longrightarrow x, y \text{ liés à } \sigma\xi, \sigma\eta
\]
\note{les lettres $\xi$, $\eta$ de ces trois lignes sont d'une lecture
incertaine, et la dernière est récrite sur une autre}
\drawing{à gauche, deux figures en perspective : un prisme de sommets $x$,
$y$, $s$, $t$ en bas et $\xi$, $\xi'$, $\eta$, $\eta'$ en haut ; au-dessous,
une seconde figure, plus grande, sur les mêmes lettres, en partie couverte de
hachures}

Prenons d'abord le lemme 6. Comme dans lemme 5, complétons les triangles
\[
(x, s, \xi), \ (y, t, \xi') \quad (\xi, \xi' \in E(w)^{*} \text{ sont liés})
\]
\[
(x, t, \eta), \ (y, s, \eta') \quad (\eta, \eta' \in E(w)^{*} \text{ sont
liés})
\]
Alors $\xi$ lié à $s$, \emph{non à $t$} (sinon $(\xi, t)$ appartiendrait à 2
triangles de sommets $x$, $\xi'$) de $=$ $\eta$ lié à $t$, non à $s$, donc cela
montre que
\[
\lbrace \pi_w(\lbrace \xi, \eta \rbrace), \complement_{P_w} \pi_w(\lbrace \xi,
\eta \rbrace) \rbrace = c_{wv}\bigl(\pi_v\lbrace s, t \rbrace, \complement
\pi_v(\lbrace s, t \rbrace)\bigr),
\]
Mais aussi $\xi$ lié à $x$, non à $y$, $\eta$ lié à $x$, non à $y$, donc
\[
\lbrace \pi_w(\lbrace \xi, \eta \rbrace), \complement \pi_w(\lbrace \xi, \eta
\rbrace) \rbrace = c_{wu}\bigl(\pi_u(\lbrace x, y \rbrace), \complement
\pi_u(\lbrace x, y \rbrace)\bigr)
\]

\page{7}
\note{« 6 », cerclé, en haut au centre}

cqfd.

Ceci étant soient $p \neq q \in E(w)^{*}$ liés à $s, t$, on sait \add{donc}
qu'ils sont non liés à $\xi, \xi', \eta, \eta'$ et non liés entre eux i.e.\
$E(w)^{*} = \lbrace \xi, \xi', \eta, \eta', p, p', q, q' \rbrace$, où $p' =
\sigma p$, $q' = \sigma q$. On a $p$ non lié à $x$ (sinon $p\,x$
appartiendrait à 2 triangles de sommets $s, t$) de $=$ $p$ non lié à $y$, donc
$p'$ lié à $x$ et $y$. de $=$ $q'$ lié à $x$ et $y$, i.e.
\[
\lbrace p', q' \rbrace = \tilde{\varphi}_{wu}(\lbrace x, y \rbrace)
\]
ce qui prouve le corollaire.
\drawing{en marge gauche, un quadrilatère $x$, $s$, $y$, $t$ (côtés $xs$,
$sy$, $yt$, $tx$) et un point $p$ au-dessus, relié à $s$ et à $t$}

\note{un trait horizontal sépare ce qui suit}

On identifie entre eux les $\Pi_{22}(P_u, P_v, P_w)$, soit $\sqsupset$ de
cardinal 3.\note{le signe qu'il donne à cet ensemble à trois éléments a la
forme d'un crochet $\sqsupset$}
[Soit $V$ le vectoriel de dim 2 sur $\mathbb{F}_2$ associé ($\sqsupset =
V^{*}$) donc $P_u, P_v, P_w$ s'identifient comme des torseurs sous $V^{*}$.
Soit $I = \lbrace u, v, w \rbrace$.

\struck{On a} Soit $\tilde{\Gamma}$ l'ens.\ des \emph{triangles} $\subset
S^{*}$, on voit que
\[
\tilde{\Gamma} \xrightarrow{\ \sim\ } \prod_{x \in E(u)^{*}}
\omega_v(x)
\]
[$E(u)^{*}$ : card.\ 8 ; $\omega_v(x)$ : card 4] donc $\operatorname{card}
\tilde{\Gamma} = 32$.\note{le signe est bien un $\prod$ ; le
cardinal $32 = 8 \cdot 4$ est celui de la somme disjointe}
Soit $\Gamma = \tilde{\Gamma}/\sigma$ (de cardinal 16)

\struck{Th} \emph{Lemme 7.} $\Gamma \to P_u \times P_v \times P_w$
[cardinal $64 = 4 \cdot 4 \cdot 4$] est injectif, et \struck{\ill{}}
\add{l'image de $\Gamma$} est un sous-torseur sous $\Psi(V) \subset V \times
V \times V$, ($\Psi(V) = \lbrace (x, y, z) \in V \times V \times V \mid x + y + z
= 0 \rbrace$).

\emph{Dém.} Soient $(x, s, \xi)$, $(y, t, \eta)$ deux triangles ayant même
image dans $P_u \times P_v \times P_w$. Quitte à appliquer $\sigma$ une ou 2
fois, \uncertain{OPS} $y = x$. Alors $s$ et $t \in E(v)^{*}$ liés à $x$ : $x$ ne
peut \struck{être} \ill{} $t = \sigma s$, donc $t = s$, de $=$

\page{8}
\note{« 7 », cerclé, en haut au centre}

$\xi = \eta$, d'où l'injectivité. Montrons que $\Gamma$ est stable par les
\uncertain{éléments} de $N_u = \operatorname{Ker}(\Psi \xrightarrow{\mathrm{pr}_1} V)$
(idem \ill{} pour $N_v$ et $N_w$) --- il s'ensuivra qu'il est stable par $\Psi
= N_u + N_v$, donc (étant de cardinal $16 = \operatorname{card} \Psi$) c'est un
torseur sous $\Psi$. Ceci équivaut à ---

\emph{Corollaire.} Soient $x \in E(u)^{*}$, $s \in E(v)^{*}$, $\xi \in
E(w)^{*}$ formant triangle. \struck{Alors} Soit \struck{\ill{}} $t \in E(v)^{*}$
\add{$t \neq s$} non lié à $s$, et lié à $\xi$ (il y en a trois), et soit $\eta$
l'élément de $E(w)^{*}$ distinct de $\xi$ lié à $s, t$ (lemme 4). Alors
\[
(\pi_u(x), \pi_v(t), \pi_w(\eta)) \in \operatorname{Im}(\Gamma \to P_u \times
P_v \times P_w),
\]
\struck{i.e.\ $\lbrace x, y, t \rbrace$ est un triangle, \ill{}} \struck{Alors}
plus précisément $(x, \sigma t, \sigma\eta)$ est un triangle.
\drawing{en marge gauche, un triangle $u$, $v$, $w$ ; plus bas, un
quadrilatère $\xi$, $s$, $\eta$, $t$ avec la diagonale $s t$, et $x$ relié à
$\xi$ par un chevron}

\emph{Dém.} Comme $x$ est lié à $s$ et $\xi$, liés resp.\ à $\eta$ et $t$, il ne
peut être lié à $t$ ni $\eta$, donc il est lié à $\sigma t$ et $\sigma\eta$, qui
sont liés entre eux (\struck{qui} puisque $t$ et $\eta$ le sont) cqfd.

\emph{NB} Les \struck{données} de $V$ et des torseurs $\Gamma$ sous $\Psi(V)$
permet de reconstituer $P_u, P_v, P_w$ et le système de « triangles » dans $P_u
\times P_v \times P_w$. Il faut \add{encore} reconstituer \struck{les} rev.\ à 2
feuillets $E(u)^{*}, E(v)^{*}, E(w)^{*}$ \add{($= \tilde{P}_u$, $= \tilde{P}_v$,
$= \tilde{P}_w$)} de $P_u, P_v, P_w$, et le revêtement à 2 feuillets
$\tilde{\Gamma}$, et les liens $\tilde{\Gamma} \to \tilde{P}_u$ etc.\ sur $\Gamma
\to P_u$ etc.]

\ldots

\page{9}
\note{« 8 », encadré, en haut au centre}

\emph{Épilogue.} Fixons un triangle $(\alpha_4, \beta_4, \gamma_4) \in
\tilde{\Gamma}$ ($\alpha_4 \in E(u)^{*}$, $\beta_4 \in E(v)^{*}$, $\gamma_4 \in
E(w)^{*}$). Alors $\Pi_{22}(P_u)$ s'identifie à $P_u - \lbrace
\pi_u(\alpha_4) \rbrace$, de $=$ \ldots. Ordonnons $V^{*}$, donc les $P_u -
\lbrace \pi_u(\alpha_4) \rbrace$ etc. Dans $E(u)^{*} - \lbrace \alpha_4
\rbrace$ sur $P_u - \lbrace \pi_u(\alpha_4) \rbrace$ il y a une section
$\lbrace \alpha_1, \alpha_2, \alpha_3 \rbrace$, définie par
\[
\mathrm{I}_0 \left\lbrace
\begin{array}{lll}
\alpha_1, \alpha_2, \alpha_3 \ (\text{et } \alpha_4) & \text{liés à} &
\beta_4 \\
\beta_1, \beta_2, \beta_3 \ (\text{et } \beta_4) & \text{---} & \gamma_4 \\
\gamma_1, \gamma_2, \gamma_3 \ (\text{et } \gamma_4) & \text{---} & \alpha_4
\end{array}
\right.
\quad ; \text{ de } = \ldots
\]
\note{les tirets de ces tableaux sont ses signes de répétition (« liés à »)}

Ainsi, en termes d'un ordre total sur $\sqsupset = V^{*}$ et d'un ordre total
(ou seulement un ordre circulaire) sur \struck{\ill{}} $I = \lbrace u, v, w
\rbrace$, on a étiqueté les $24$ \add{$= 2 \times 12 = 2 \times (3 \times
4)$} élts.\ de $S^{*}$ (les pts autres que $u, v, w$) par les
\[
\alpha_i, \ \beta_i, \ \gamma_i, \qquad \alpha'_i = \sigma\alpha_i, \
\beta'_i = \sigma\beta_i, \ \gamma'_i = \sigma\gamma_i .
\]

De (I), et lemme 3, on déduit
\[
[0] \quad \alpha_i \text{ lié à } \alpha'_i \ / \ \beta_i \text{ à }
\beta'_i \ / \ \gamma_i \text{ à } \gamma'_i
\]
\[
\mathrm{I}' \left\lbrace
\begin{array}{llll}
\alpha'_i & \text{liés à } \beta'_4, & \text{non à } \beta_4 & \text{ni
entre eux} \\
\beta'_i & \text{---} \ \gamma'_4, & \text{non à } \gamma_4 & \text{ni entre
eux} \\
\gamma'_i & \text{---} \ \alpha'_4, & \text{non à } \alpha_4 & \text{ni entre
eux}
\end{array}
\right.
\qquad
\mathrm{I} \left\lbrace
\begin{array}{llll}
\alpha_i & \text{liés à } \beta_4, & \text{non à } \beta'_4 & \text{ni entre
eux} \\
\beta_i & \text{---} \ \gamma_4 & \text{---} \ \gamma'_4 & \text{---} \\
\gamma_i & \text{---} \ \alpha_4 & \text{---} \ \alpha'_4 & \text{---}
\end{array}
\right.
\]
\note{sur la page, (I′) est écrit au-dessus de (I) et une flèche les relie ;
dans la première ligne de (I′) quelques lettres après $\alpha'_i$ sont
biffées. La première ligne de (I′) porte $\beta_4$, $\beta_4$ là où le
tableau (I) invite à lire $\beta'_4$, $\beta_4$ : les accents sont d'une
lecture incertaine}

\emph{Lemme 8} ($1 \leqslant i \leqslant 3$)
\[
\mathrm{II} \left\lbrace
\begin{array}{lll}
\alpha_i \text{ liés à } \gamma'_4, & \text{non à } \gamma_4 \\
\beta_i \ \text{---} \ \alpha'_4, & \text{---} \ \alpha_4 \\
\gamma_i \ \text{---} \ \beta'_4, & \text{---} \ \beta_4
\end{array}
\right.
\qquad
\mathrm{II}' \left\lbrace
\begin{array}{lll}
\alpha'_i \ \text{---} \ \gamma_4, & \text{non à } \gamma'_4 \\
\beta'_i \ \text{---} \ \alpha_4, & \text{---} \ \alpha'_4 \\
\gamma'_i \ \text{---} \ \beta_4, & \text{---} \ \beta'_4
\end{array}
\right.
\]

\page{10}
\note{« 9 », cerclé, en haut au centre}

Il suffit de prouver que $\alpha_i$ non lié à $\gamma_4$, mais comme il est lié
à $\beta_4$, $(\beta_4, \gamma_4)$ serait \struck{\ill{}} \add{côté} de deux
triangles (de sommets $\alpha_i$ et $\alpha_4$), cqfd
\drawing{à gauche, un triangle $\alpha_4$, $\beta_4$, $\gamma_4$, et un trait
de $\beta_4$ à $\alpha_i$}

\emph{Lemme 9}
\[
\mathrm{III} \left\lbrace
\begin{array}{lll}
\alpha_i \text{ lié à } \beta_j, \gamma_j, \beta'_i, \gamma'_i & (1
\leqslant i \neq j \leqslant 3) & \text{non à } \beta'_j, \gamma'_j, \beta_i,
\gamma_i \\
\beta_i \text{ lié à } \gamma_j, \alpha_j, \gamma'_i, \alpha'_i & (\text{---})
& \text{---} \ \gamma'_j, \alpha'_j, \gamma_i, \alpha_i \\
\gamma_i \ \text{---} \ \alpha_j, \beta_j, \alpha'_i, \beta'_i & (\text{---}) &
\text{---} \ \alpha'_j, \beta'_j, \alpha_i, \beta_i
\end{array}
\right.
\]
\[
\mathrm{III}' \left\lbrace
\begin{array}{lll}
\alpha'_i \text{ lié à } \beta'_j, \gamma'_j & \text{--} & \text{--} \\
\text{--} & \text{--} & \text{--} \\
\text{--} & \text{--} & \text{--}
\end{array}
\right.
\]
\note{dans (III), les deux derniers termes des colonnes « lié à » sont récrits
sur d'autres lettres et leur accent est d'une lecture incertaine ; (III′)
n'est écrit que pour sa première ligne, le reste en signes de répétition. La
première lettre de (III′) est lue $\alpha'_i$ par symétrie avec (II′) ; la
page porte peut-être $\alpha_i$}

Il reste à voir que $\alpha_1$ lié à $\beta_2$, \struck{$\beta_3$} \add{non à
$\beta_1$}. Or
\[
\begin{array}{ll}
\alpha_4 \text{ lié à} & \beta'_1, \beta'_2, \beta'_3, \beta_4 \quad
(\text{lemme } 8) \\
\alpha_1 \text{ lié à} & ? \quad ? \quad ? \quad \beta_4
\end{array}
\]
\struck{Donc} l'ens.\ des coïncidences entre $\omega_v(\alpha_1)$ et
$\omega_v(\alpha_4)$ doit être dans la partition $(1, 4)\,(2, 3)$, se donc
être \uncertain{$\beta_4$} $(1, 4)$ puisqu'il contient 4, donc
\[
\alpha_1 \text{ lié à } \beta'_1, \beta_2, \beta_3, \beta_4 \qquad
\text{cqfd.}
\]
\note{« se donc être » est lu tel quel ; la lettre devant $(1,4)$ est chargée
d'encre}

\emph{Conséquence :} $S$ est déterminé par les conditions de départ, à isom.\
près. \struck{\ill{}} donc un graphe cubique

cqfd

\note{en bas à droite : « TSVP »}

\page{11}
\note{« 10 », cerclé, en haut au centre}

De plus, on trouve ceci :

Considérons les \struck{diagrammes} \add{graphes} suivants
\[
\begin{array}{lll|l}
D_0 = \emptyset & & & 1 \\
\cap \ D_1 = \bullet\, u & \text{pt} & & 27 = 3^3 . \\
D_2 = u \,\text{---}\, v & \text{arête} & & 27 \cdot 10 = 2 \cdot 3^3 \cdot 5 = 270 \\
\cap \ D_3 = \text{triangle } u, v, w & \text{triangle} & & 27 \cdot 10 \cdot 1 = 2 \cdot 3^3 \cdot 5 = 270 \\
D_4 & \text{\struck{triangle} \add{balise}} & & 27 \cdot 10 \cdot 1 \cdot 8 = 2^4 \cdot 3^3 \cdot 5 = 2160 \\
\cap \ D_5 & \text{préprisme (triangulaire)} & & 27 \cdot 10 \cdot 1 \cdot 8 \cdot 4 = 2^6 \cdot 3^3 \cdot 5 = 8640 \\
\cap \ D_6 & \text{prisme (triangulaire)} & & 27 \cdot 10 \cdot 1 \cdot 8 \cdot 4 \cdot 1 = 2^6 \cdot 3^3 \cdot 5 = 8640 \\
\cap \ D_7 & \text{bitriangle} & & 27 \cdot 10 \cdot 1 \cdot 8 \cdot 4 \cdot 1 \cdot 1 = 2^6 \cdot 3^3 \cdot 5 = 8640
\end{array}
\]
\note{la colonne de gauche est une suite de dessins, reliés par des signes
$\cap$ d'inclusion ; $D_1$, $D_3$ et $D_7$ sont cerclés. La colonne de droite
est séparée des autres par un trait vertical}
\drawing{les graphes $D_i$ : $D_1$ le point $u$ ; $D_2$ l'arête $u v$ ;
$D_3$ le triangle $u, v, w$ ; $D_4$ le triangle $u, v, w$ avec une arête
pendante $u u'$ ; $D_5$ le triangle $u, v, w$ et, au-dessous, $u'$, $v'$
(la moitié inférieure d'un rectangle hachurée, reprise) ; $D_6$ un prisme
triangulaire à bases $u, v, w$ et $u', v', w'$ ; $D_7$ deux prismes accolés
par une face, avec les sommets $u, v, w, u', v', w', w''$ et $w'''$, dessin
repris sur une première esquisse hachurée}

Soit \add{$\mathrm{Pl}_i =$} $\mathrm{Pl}(D_i, C)$ l'ens.\ des
\emph{plongements} \add{(« pleins »)} de $D_i$ dans $C$, $W = \operatorname{Aut}
C$, $G_i = \operatorname{Aut} D_i$. Alors\note{sa lettre $W$ est barrée d'un
trait oblique}

\struck{b)} \add{c)} Les $\mathrm{Pl}(D_i, C)$ sont des espaces
\emph{homogènes} (non vides) sous $W$

a) Les morphismes $\mathrm{Pl}(D_{i+1}, C) \to \mathrm{Pl}(D_i, C)$ ($0
\leqslant i \leqslant 6$) transposés des inclusions $D_i \subset D_{i+1}$ sont
\emph{surjectifs} \uncertain{et} sont des rev.\ d'ordres respectifs $27, 10, 1,
8, 4, 1, 1$

b) (Par conséquent) les cardinaux des $\mathrm{Pl}_i$ sont les nbs marqués
dans la colonne de droite.)\note{une flèche fait passer le c) avant le a)}
\struck{Donc} \add{$\overline{C}$} Donc les cardinaux \struck{des} de
l'ens.\ \add{$\operatorname{Conf}_{D_i}(C)$} des sous-graphes \add{pleins}
\struck{$\operatorname{Conf}_{D_i}(C)$} de $C$ de type $D_i$ (s'identifiant à
$\mathrm{Pl}_i/G_i$) est donné par $\operatorname{card}
\mathrm{Pl}_i/\operatorname{card} G_i$, qui est
\[
\begin{array}{llll}
\operatorname{Conf}_{D_0} & 1 & \text{---} \operatorname{Aut} D_0 & e \\
\operatorname{Conf}_{D_1} & 27 & \text{---} \operatorname{Aut} D_1 & e \\
\operatorname{Conf}_{D_2} & 27 \cdot 10 : 2 = 135 & \text{---}
\operatorname{Aut} D_2 & \mathfrak{S}_2 \\
\operatorname{Conf}_{D_3} & 27 \cdot 10 \cdot 1 : 6 = 45 & \text{---}
\operatorname{Aut} D_3 & \mathfrak{S}_3 \\
\operatorname{Conf}_{D_4} & 27 \cdot 10 \cdot 1 \cdot 8 : 2 = 2^3 \cdot 3^3
\cdot 5 = 1080 & \text{---} \operatorname{Aut} D_4 & \mathfrak{S}_2 \\
\operatorname{Conf}_{D_5} & 27 \cdot 10 \cdot 1 \cdot 8 \cdot 4 : 2 = 2^5
\cdot 3^3 \cdot 5 = 4320 & \text{---} \operatorname{Aut} D_5 &
\mathfrak{S}_2 \\
\operatorname{Conf}_{D_6} & 27 \cdot 10 \cdot 1 \cdot 8 \cdot 4 \cdot 1 : 12
= 2^4 \cdot 3^2 \cdot 5 = 720 & \text{---} \operatorname{Aut} D_6 &
\mathfrak{S}_2 \times \mathfrak{S}_3 \\
\operatorname{Conf}_{D_7} & 27 \cdot 10 \cdot 1 \cdot 8 \cdot 4 \cdot 1
\cdot 1 : 72 = 2^3 \cdot 3 \cdot 5 = 120 & \text{---} \operatorname{Aut} D_7 &
\mathfrak{S}_2 . (\mathfrak{S}_3 \times \mathfrak{S}_3)
\end{array}
\]
\note{dans la dernière ligne, « 72 » et l'exposant de 3 sont récrits, et un
premier groupe est biffé avant $\mathfrak{S}_2 . (\mathfrak{S}_3 \times
\mathfrak{S}_3)$. $C$ désigne ici, semble-t-il, le graphe $S$ des pages
précédentes}

\page{12}
\note{« 10 bis », cerclé, en haut au centre}

d) Le stabilisateur d'un élément de $\mathrm{Pl}_6$ \struck{dans} \add{(ou, ce
qui revient au $=$, de $\mathrm{Pl}_5$ ou $\mathrm{Pl}_7$)} dans $W$ est
isomorphe à $\mathfrak{S}_3$ (via l'opération sur l'ens.\ $I$ \add{de cardinal
3} des éléments de \struck{$E(u)^{*} = E(u) - \lbrace v, w \rbrace$} $E(u)^{*} =
E(u) - \lbrace v, w \rbrace$ qui sont $\neq u'$, \struck{non liés à $v$} et liés
à $v'$, i.e.\ des él.\ de $E(u) - \lbrace v, \struck{w}\, u' \rbrace$ qui sont
liés à $v'$, i.e.\ $I = E(u) \cap E(v') - \lbrace u', v \rbrace$. On voit
\add{donc} que cet $I$ est can.\ is.\ aux $(9 \cdot 4) : 2 = 18$ ensembles
analogues tels que $E(v'') \cap E(w') - \lbrace u'', w' \rbrace$ etc
\ldots\ $\rbrace$. Donc $W$ est de cardinal égal à : $\operatorname{card}
\mathrm{Pl}_6 \cdot 6$, i.e.
\[
\operatorname{card} W = 2^7 \cdot 3^4 \cdot 5 = 51840
\]
\note{un double trait sépare ce qui suit}

\struck{\uncertain{Complexe} cubique} \add{\uncertain{Compte}}\note{deux
mots biffés au bas de la page, un mot récrit au-dessus ; la page s'arrête là}

\page{13}
\note{« 11 », cerclé, en haut au centre}

La catégorie formée des $S$, avec un triangle $I \subset S$ et un triangle
$\Pi' \subset S$ disjoint du précédent (pour les iso) équivaut à celle des
couples $(I, \Pi)$ de deux ens.\ de cardinal 3.\note{les deux ensembles sont
écrits $I$ et $\Pi$, ce dernier comme une potence ; le « $\Pi'$ » de la
deuxième ligne est d'une lecture incertaine}

À un \add{tel} couple, on associe
\[
S = I \amalg (I \amalg I) \amalg I \times \Pi \times \omega(I)
= (I \times \lbrace 0, 1, 2 \rbrace) \amalg (I \times \Pi \times \omega(I)),
\]
\[
I \ni u \longmapsto \alpha(u), \ \alpha'(u)
\]
où $\omega(I)$ \add{$\subset \mathfrak{S}_I$} est l'ens.\ \add{(de cardinal
2)} des ordres circulaires sur $I$.

\struck{Si $i \in I$, $i$ est lié à}
\drawing{biffé, à gauche : $I$ et trois flèches parallèles marquées $i_3$,
$i'_1$, $i_2$}

\[
\begin{array}{lll}
(i, \nu) \text{ et } (j, \mu) \ (\text{distincts}) & \text{sont liés ssi} &
i = j \text{ \emph{ou} } \nu = \mu \\
(i, \nu) \text{ et } (j, \pi, \rho) & \text{liés dans les cas} &
\left\lbrace
\begin{array}{ll}
i = j : & \nu = 0 \\
i \neq j & \left\lbrace \begin{array}{l} \nu = 1, \ j = \rho i \\ \nu =
2, \ i = \rho j \end{array} \right.
\end{array}
\right. \\
(j, \pi, \rho) \text{ et } (j', \pi', \rho') \ \text{distincts} & \text{liés
dans les cas} &
\left\lbrace
\begin{array}{ll}
j = j' & \pi = \pi', \ (\rho \neq \rho') \\
j \neq j' & \left\lbrace \begin{array}{ll} i \neq j, & \rho = \rho' \\ i =
j & \rho \neq \rho' \end{array} \right.
\end{array}
\right.
\end{array}
\]
($j \in I$, $\pi \in \Pi$, $\rho \in \omega(I)$)\note{dans le deuxième cas
une première rédaction est raturée ; dans le troisième, les $i$ sont écrits
tels quels, bien que seuls $j$, $j'$ figurent dans le couple}
\drawing{à gauche, un prisme vu d'en haut : triangle $u, v, w$ en bas,
sommets $\alpha_4$, $\beta_4$, $\gamma_4$ au milieu (le triangle $\alpha_4
\gamma_4 \beta_4$ repassé), $\alpha'_4$, $\beta'_4$, $\gamma'_4$ à
l'extérieur}

L'automorphisme $\sigma$ conserve $I \times \lbrace 0, 1, 2 \rbrace$ et son
complémentaire, et agit sur le premier via l'identité sur $I$, et la
permutation entre 1 et 2 sur l'autre facteur ; sur $I \times \Pi \times
\omega(I)$ il agit par permutation des deux \struck{\ill{}} élts de $\omega(I)$
(et l'identité sur $I \times \Pi$).

\marginal{à gauche : « $u$, $\alpha_4$, $\alpha'_4$ » et, dessous,
« $(\ldots_i)_{1 \leqslant i \leqslant 3}$ » (raturé), « $\alpha_i,
\alpha'_i$ », « $\beta_j, \beta'_j$ », « $1 \leqslant i, j \leqslant 3$ »,
« $\alpha_4 \mapsto \alpha'_4$ »}

\[
I \amalg (I \times \varepsilon) \amalg (I \times \Pi \times \varepsilon
\times \omega(I)/\mathbb{F}_2)
\]
\[
(I \times J) \amalg
\]
\drawing{en bas à gauche, une grille de droites croisées dont les cases
portent $\alpha_4$, $\beta_4$, $\gamma_4$, $\alpha_1$, $\beta_2$,
$\gamma_3$ et d'autres lettres, en partie illisibles}

\page{14}
\note{« 12 », cerclé, en haut au centre}

\emph{Proposition.} $G$ groupe abélien (pour simplifier), $V = G \otimes_{\mathbb{Z}}
\mathbb{F}_2$, $\varepsilon$ un $\mathbb{F}_2$-vectoriel, $\Gamma = \Psi(G) =
\operatorname{Ker}(G \times G \times G \to G)$, $(x, y, z) \mapsto x + y + z$,
$p_i : \Psi(G) \to G$ induits par $\mathrm{pr}_i$ ($1 \leqslant i \leqslant 3$),
$N_i = \operatorname{Ker} p_i \rightleftarrows G$ (via $\alpha_i$) \struck{$\lbrace
(x, \hat{x}, x \rbrace$}.\note{la lettre qu'il donne à $\Psi(G)$ est un $\Gamma$
repassé ; la flèche double porte $\alpha_i$ en dessous}

a) Pour toute extension centrale
\struck{$E_0$} :
\[
1 \to \varepsilon \to E_0 \to G \to 1
\]
considérons $\psi(E_0) = p_1^{*}(E_0) \wedge p_2^{*}(E_0) \wedge p_3^{*}(E_0) =
\tilde{\Gamma}$ (produit de Baer d'extensions de \struck{\ill{}} $\Gamma$ par
$\varepsilon$, image inverse de l'ext.\ $E_0$ (de $G$ par $\varepsilon$) par les
$p_i$). Comme
\[
p_j \alpha_i = \left\lbrace
\begin{array}{ll}
0 & \text{si } i = j \\
\mathrm{id}_G & \text{si } i \neq j
\end{array}
\right. ,
\]
on trouve que $\alpha_i^{*}(\tilde{E})$ est splitté canoniquement.\note{« produit
de Baer » : lecture incertaine du mot qui suit « produit »}
d'où

\begin{tikzcd}
1 \arrow[r] & \varepsilon \arrow[r] & \tilde{E} \arrow[r] & \Gamma \arrow[r] & 1 \\
 & & & G \arrow[ul, "\rho_i"] \arrow[u, "\alpha_i"'] &
\end{tikzcd}

Ceci posé, soit
\[
q : \dot{V} \to \varepsilon
\]
\uncertain{est} l'application quadratique qui définit la classe de l'extension
$E_0$, on a pour $g \in G$
\[
\rho_1(g)\, \rho_2(g)\, \rho_3(g) = q(\dot{g})
\]
où $\dot{g} \in V$ est la classe de $g \in G$ dans $V = G/2G$.\note{dans cette formule, un $\alpha_1$, $\alpha_2$,
$\alpha_3$ biffé suit chaque $\rho_i$ ; le point au-dessus de $V$ dans la
ligne de $q$ est net sur la page}

b) Toute \add{syst.\ $(\tilde{E}, \rho_1, \rho_2, \rho_3)$ d'une} extension
\add{centrale} $\tilde{E}$ de $\Gamma$ par $\varepsilon$ et de relèvements
$\rho_1, \rho_2, \rho_3$ de $\alpha_1, \alpha_2, \alpha_3$ provient d'une
extension $E_0$ de $G$ par $\varepsilon$, déterminée : isom.\ (\uncertain{à un}
unique) près --- sa classe $q \in \operatorname{Quad}_{\mathbb{Z}}(G, \varepsilon)
\simeq \operatorname{Quad}_{\mathbb{F}_2}(V, \varepsilon)$ est donnée par la
formule précédente. La catégorie \add{(pour les iso)} des \struck{triples}
syst.\ $(\tilde{E}, \rho_1, \rho_2, \rho_3)$ est \emph{rigide} i.e.\ tout
automorphisme est l'identité (donc \uncertain{tout automorphisme} de l'ext.\ $E_0$ induit
l'identité dans \add{l'ext} $\psi(E_0) \ldots$)

\page{15}
\note{« 13 », cerclé, en haut au centre}

c) Conditions équivalentes \add{$2G = 0$ i.e.\ $G \simeq V$, et} \struck{\ill{}}

(i) \add{$q(x) \neq 0$ pour} \struck{\ill{}} $x \in V^{*} = V - \lbrace 0
\rbrace$

(ii) Posant $\tilde{P}_i = \tilde{E}/\rho_i(G)$, \struck{la fam Posons} et
$\tilde{p}_i : \tilde{E} \to \tilde{P}_i$, si $x_1, x_2, x_3 \in \tilde{\Gamma}$
sont tels que
\[
p_1(\tilde{x}_2) = p_1(\tilde{x}_3), \quad p_2(\tilde{x}_3) =
p_2(\tilde{x}_1), \quad p_3(\tilde{x}_1) = p_3(\tilde{x}_2),
\]
alors $\tilde{x}_1 = \tilde{x}_2 = \tilde{x}_3$.\note{les indices de ces
égalités sont d'une lecture incertaine, surtout ceux de la dernière ligne}

(iii) Si $\tilde{z}_1 \in \rho_1(G)$, $\tilde{z}_2 \in \rho_2(G)$, $\tilde{z}_3
\in \rho_3(G)$ sont tels que $\tilde{z}_1 \tilde{z}_2 \tilde{z}_3 = 0$, alors
$\tilde{z}_1 = \tilde{z}_2 = \tilde{z}_3 = 1$

(iv) Si $g \in G$ est tel que $\rho_1(g) \rho_2(g) \rho_3(g) = 1$, alors $g =
1$\note{le signe $=$ est repassé}

De plus, pour $G$, $\varepsilon$ fixés avec $\varepsilon$ \uncertain{$\neq$}
$\mathbb{F}_2$, $\exists$ une $\tilde{\Gamma}$ \struck{est} avec ces propriétés
ssi $2G = 0$ i.e.\ $G \simeq V$, $\operatorname{rg}_{\mathbb{F}_2} V \leqslant
2$ (et cette extension est unique à isom.\ unique près).

\page{16}
\note{« 13 bis », cerclé, en haut au centre}

$S$ graphe \add{$\neq \emptyset$} tel que \struck{l'ens.\ des \ill{} \ill{}
$\neq \emptyset$ et}

a) $\forall x, y \in S$ liés, $\exists!\, z \in S$ tel que $z$ soit lié à $x$ et
$y$

b) $\exists\, \nu \in \mathbb{N}^{*}$ tel que $\forall s \in S$,
$\operatorname{card} E(s) = \nu$ \quad ($E(s) = \lbrace t \in S \mid t \text{
lié à } s \rbrace$)

\struck{c)} \emph{Lemme 1} \add{$\nu$ est pair, et} On a $\operatorname{card} S
\geqslant 3(\nu - 1)$. On a égalité ssi $\forall$ \add{triangle (\uncertain{il
en} $\exists$)} $T = \lbrace x, y, z \rbrace$ de $S$, \struck{a} tt élément de
$S^{*} = S - T$ est lié à un (et un seul) élément de $T$.\note{l'inégalité
porte $3(\nu - 1)$ sur la page, récrit sur un autre chiffre ; le théorème de
la page 2 ($\nu = 10$, $\operatorname{card} S = 27$) correspond à
$3(\nu - 1)$}

\[
\nu = 2(\mu + 1), \qquad \operatorname{card} S = 3(2\mu + 1)
\]
\[
\operatorname{card} S - E(s) - \lbrace s \rbrace \ \ldots \
\operatorname{card} (S - E(s) - \lbrace s \rbrace) = 4\mu
\]
\[
4\mu \leqslant 2^{\mu + 1}, \qquad \mu \leqslant 2^{\mu - 1}
\]
\note{calculs griffonnés, avec des surcharges ; un « 2 » isolé au-dessus de
la deuxième ligne et un mot raturé au-dessous}
\drawing{un petit triangle ; à côté, un graphe à quelques sommets entouré de
deux ovales qui se croisent ; plus bas, une étoile à trois branches dont
chaque bout porte un petit triangle (deux sommets reliés au bout de la
branche)}

\page{17}
\note{« 14 », cerclé, en haut au centre}

\emph{Graphes cubiques \add{$C$} : sommet marqué $u$}

Revient à la donnée d'un ens.\ $Q_u$ de cardinal 5 (l'ens.\ des triangles de
sommet $u$), i.e.\ ens.\ des 5 couples de pts mutuellement liés de $E(u)$), et
d'un torseur $P$ sous $\Psi = \Psi(Q_u, \mathbb{F}_2) = \operatorname{Ker}
(\mathbb{F}_2^{Q_u} \to \mathbb{F}_2)$, $(x_i)_{i \in I} \mapsto \sum x_i$.
\struck{On reconstitue} \struck{le torseur}\note{la lettre de l'ensemble à cinq
éléments, repassée à l'encre plus foncée, tient du $P$ et du $Q$ cursifs ; la
page suivante l'écrit nettement $Q_u$, et elle est rendue ainsi}
$E(u)$ comme l'ens.\ somme des $P \wedge_{\Psi} (\mathbb{F}_2, p_i)$, où les
$p_i : \Psi \to \mathbb{F}_2$ \add{($i \in E$)} sont induits par les projections
$\mathbb{F}_2^{Q_u} \xrightarrow{\mathrm{pr}_i} \mathbb{F}_2$ ($i \in Q_u$), et
$C - E(u)$ comme $P$, donc
\[
C \simeq \Bigl( \coprod_{i \in E} P \wedge_{\Psi} (\mathbb{F}_2, p_i) \Bigr)
\amalg P \amalg \lbrace u \rbrace .
\]
La recette pour ce dictionnaire est la suivante : $E(u)$ est un ens.\ de
cardinal 10 muni d'un automorphisme involutif \uncertain{d'ordre} sans pt fixe
(car tt él.\ de $E(u)$ est lié à un et un seul él.\ de $E(u)$, grâce à
l'axiome des triangles), or une telle structure est déterminée par un $Q_u$ de
cardinal 5 ($Q_u = E(u)/\sigma$) et une famille d'ens.\ à 2 él.\ i.e.\ de
$\mathbb{F}_2$-torseurs, paramétrisés par $Q_u$ --- i.e.\ un torseur $Q$ sous
$(\mathbb{F}_2)^{Q_u}$. D'autre part, tt $x \in C - E(u)$ définit une partie de
$E(u)$ --- savoir $E(x) \cap E(u)$ --- telle que sa représentation avec les
classes sous $\lbrace e, \sigma \rbrace$ i.e.\ les fibres de $E(u) \to Q_u$
\add{$= \tilde{Q}_u$} i.e.\ les orbites $\subset E(u)$, soit de cardinal 1. Or
l'ens.\ de ces parties n'est autre que $Q$ (qui est de cardinal $2^5 = 32$)
alors que $C - E(u)$ est de card $16 = 32/2$. On trouve que $C - E(u) \to Q$ est
\emph{injectif}, et son image \add{$P$} est une torseur sous $\Psi \subset
\mathbb{F}_2^{Q_u}$ ($\Psi = \operatorname{Ker}(\mathbb{F}_2^{Q_u} \to
\mathbb{F}_2)$) --- donc on a sur $Q$ une restriction du groupe\note{« $C -
E(u)$ » s'entend ici, semble-t-il, privé de $u$ : $1 + 10 + 16 = 27$}

\page{18}
\note{« 15 », cerclé, en haut au centre}

structural : $\Psi$. Donc la donnée de \add{$Q_u$,} $Q$ avec cette restriction
équivaut à celle de $Q_u$, $P$. Pour reconstituer le graphe $C$, il faut
reconstituer la notion de « sommets liés » sur
\[
\lbrace u \rbrace \cup \underbrace{\Bigl( \coprod_i P \wedge_{\Psi}
(\mathbb{F}_2, p_i) \Bigr)}_{E_p^{*}} \amalg P .
\]
\note{sous l'accolade, après $E_p^{*}$, un signe raturé}
On déclare :

a) $u$ est lié aux pts de $E_p^{*}$, et non à ceux de $P$

b) \struck{un} \add{a} pt $a$ de $E_p^{*}$ est lié à $\sigma a$ [i.e.\ si $a
\in P \wedge_{\Psi} (\mathbb{F}_2, p_i)$, il est lié à l'autre pt $a'$ de cet
ens.] et à aucun autre pt de $E_p^{*}$

c) un pt $a$ de $E_p^{*}$ est lié à un pt \struck{\ill{}} $x$ de $P$ ssi,
prenant $a \in P \wedge_{\Psi} (\mathbb{F}_2, p_i)$, $a$ est dans l'image de
$x$ par l'application canonique $P \to P \wedge_{\Psi} (\mathbb{F}_2, p_i)$.

d) Deux pts $x, y$ de $P$ sont liés ssi l'élément $y - x = \alpha \in
\Psi(Q_u) = \operatorname{Ker}(\mathbb{F}_2^{Q_u} \to \mathbb{F}_2)$ est
\struck{\ill{}} de la forme \struck{\ill{}} $e'_i$ ($i \in Q_u$), où
\[
e'_i = \sum_{j \in Q_u - \lbrace i \rbrace} e_j \in \Psi
\]
($(e_j)_{j \in Q_u}$ étant la base canonique de $\mathbb{F}_2^{Q_u}$). Ceci
se fonde sur le\note{dans (d), « $\sum$ » est surmonté d'un trait ; la
lecture « se fonde sur le » est incertaine}

\emph{Lemme} Soient $u \in C$, $x, y \in C - E(u) - \lbrace u \rbrace$. Pour que
$x$ et $y$ soient liés, il faut et il suffit que \add{l'on} \struck{puisse
\ill{} de l'ens.} \add{existe} \struck{un triangle} \add{(un et) un seul}
triangle \struck{des 5 triangles} $t$ de sommet $u$ tel que
\[
t - E(x) = t - E(y)
\]
(NB les deux membres sont \add{de la forme} \struck{réduits} \add{réduit à un
él.\ $\neq u$}) \struck{réduits} de cardinal 1) OPS $x \neq y$.\note{les
reprises de cet énoncé s'enchevêtrent ; l'ordre des insertions est celui
qu'on a pu suivre}

\emph{Dém.} S'il n'existait aucun tel triangle, on aurait $x - y = \sum_{i \in
Q_u} e_i$, or $\sum e_i \notin \Psi$, \add{car $5 \neq 0$ dans $\mathbb{F}_2$},
absurde. Soit $\lbrace u, v, w \rbrace$ un tel triangle, donc $x, y$ sont liés
tous deux à $v$ ou $w$ --- disons $v$. Si $y$ et $x$ sont liés, \uncertain{la
loi} (ou la condition \uncertain{habituelle}) $y = \sigma x$, on sait que pour
tt $a \in E(u)^{*}$ \add{$= E(u) - \lbrace v, w \rbrace$}, si $x$ lié à $a$,
$\sigma x$ non lié à $a$ et inversement.

\page{19}
\note{« 16 », cerclé, en haut au centre}

donc $t = \lbrace u, v, w \rbrace$ est unique de son espèce. Si $x$ et $y$
sont non liés, on sait qu'il existe un \add{(et $=$ deux)} $a \in E(u)^{*}$
tel que $a$ lié à la fois à $x$ et $y$, donc $t$ \uncertain{n'est} unique de
son espèce, cqfd.

\emph{Corollaire} Le groupe \struck{\ill{}} $W_u$ \add{$\subset W$} des
automorphismes de $C$ qui fixent $u$ est \struck{\ill{}} isomorphe au groupe
des couples $(\varphi, \bar{\varphi})$ d'un automorphisme $\varphi \in
\mathfrak{S}_{Q_u}$ de $Q_u$, et d'un $\bar{\varphi}$ automorphisme du
\struck{\ill{}} $\Psi_{Q_u}$-torseur $P$, où $\bar{\varphi}$ est l'automorphisme
de $\Psi_{Q_u}$ induit par $\varphi$. Donc on a une suite exacte
(splittée)

\begin{tikzcd}
1 \arrow[r] & \Psi_{Q_u} \arrow[r] \arrow[d, no head, "\wr"'] & W_u \arrow[r] & \mathfrak{S}_{Q_u} \arrow[r] \arrow[d, no head, "\wr"] & 1 \\
 & \mathbb{F}_2^4 & & \mathfrak{S}_5 &
\end{tikzcd}

\note{au-dessus de $(\Psi_{Q_u})$, un mot biffé ; les deux isomorphismes
verticaux sont écrits $\wr$}

donc $W_u \simeq$ groupe de Weyl $W_{D_5}$ de $D_5$ ($\simeq SO(10)$) qui
\uncertain{apparaît} donc comme sous-groupe d'indice 27 dans $W$ ($\simeq
W_{E_6}$).

\emph{NB} Il faudrait expliciter comment la donnée de $(C, u)$ permet de
définir un syst.\ de racines $D_5$ (plongé dans le syst.\ de racines $E_6$
associé à $C$) --- donc définir un ens.\ de \add{40} « bases » (\uncertain{=
racines}) de $C$, stable par $W$ : bien précises. Or on trouve qu'il y a $40 =
(10 \cdot 8) : 2$ \add{$= \binom{10}{2} - 5$} \struck{couples} paires
d'éléments de $E(u)$ non liés. Or \ill{} le

\emph{Lemme} Pour tt paire \add{$\lbrace \xi_1, \xi_2 \rbrace$} d'éléments non
liés de

\page{20}
\note{« 17 », cerclé, en haut au centre}

$E(u)$, $\exists!$ base $b \subset C$ telle que $b \cap E(u) = \lbrace \xi_1,
\xi_2 \rbrace$ et (si $b'$ est la base opposée) on a $b' \cap E(u) = \lbrace
\sigma\xi_1, \sigma\xi_2 \rbrace$, \struck{\ill{}} où $\sigma$ est
l'automorphisme can.\ d'ordre 2 de $E(u)$.

\add{NB On aura donc $\sigma\xi_1 = \xi'_2$, $\sigma\xi_2 = \xi'_1$. Alors $b
\cap E(u) = \lbrace \xi_1, \xi_2 \rbrace$, $b' \cap E(u) = \lbrace \sigma\xi_1,
\sigma\xi_2 \rbrace$.}\note{ligne ajoutée entre les lignes et rattachée par un
trait ; les accents des $\xi'$ sont d'une lecture incertaine}

\marginal{en oblique dans la marge gauche : « \uncertain{rappeler} \uncertain{plus
simple} \uncertain{plus bas} »}
\emph{Dém.} Soit une base $\lbrace \xi_1, \xi_2, \ldots, \xi_6 \rbrace$
satisfaisant \uncertain{à la \ill{}} \ldots\ On aura alors \struck{\ill{}} que
$u$ est lié à $\xi_1, \xi_2, \xi'_1, \xi'_2$, donc $u = \bar{\xi}_{12}$, et $b
\cap E(u) = \lbrace \xi_1, \xi_2 \rbrace$, $b' \cap E(u) = \lbrace \xi'_1,
\xi'_2 \rbrace$. Pour l'existence \add{d'une base $b$ satisfaisant les
conditions envisagées}, on voit que \uncertain{comme} le groupe $W_u$ est
transitif sur l'ens.\ des \struck{couples} \add{paires} $\lbrace \xi_1, \xi_2
\rbrace$ non liés de $E(u)$ et $W = W_{E_6}$ transitif sur l'ens.\ des $v \in
C$) on peut supposer, p.ex.\ : une base $(\xi_1, \ldots, \xi_6)$ déjà donnée,
que $u = \bar{\xi}_{12}$, $\xi_1 = \bar{\xi}_1$, $\xi_2 = \bar{\xi}_2$.\note{la
lettre de $W_{E_6}$ est récrite sur une autre}

Pour l'unicité, on voit que \struck{\ill{} \ill{} \ill{} $(3 \leqslant i
\leqslant 6)$ \ill{} lié à $\xi_1$, $\sigma\xi_1 = \xi'_2$, et $\sigma\xi_2 =
\xi'_1$, \ill{} \ill{} distinct de $\xi_1$, \ill{} \ill{} $E(u)$, \ill{} \ill{}
d'après ce qui \ill{} \ill{} les $\bar{\xi}_i$, $\bar{\xi}'_i$ $(3 \leqslant i
\leqslant 6)$ \ill{} $\in E(u)$, donc que \ill{} \ill{} par des
$\xi_j$, $\xi'_j$}, disons que $\bar{\xi}_3$ devrait \ill{} : $\xi_1, \xi_2$ et
lié à $\xi'_1, \xi'_2$ donc $u$ n'est pas un $\xi'_j$, ni un $\xi_j$ (car ce
serait un $\xi_j$ --- $\xi_j$ par \ill{} lié à $\xi'_1, \xi'_2$, mais il le
serait aussi à $\xi_1, \xi_2$, absurde), de $=$ pour les $\bar{\xi}'_i$ $(3
\leqslant i \leqslant 6)$ cqfd.\note{le milieu de ce paragraphe est un bloc de
lignes biffées, encadré et barré de deux grands traits en croix ; seuls
quelques signes s'y lisent}
\drawing{à gauche du bloc biffé, deux segments issus d'un même point, aux
extrémités marquées $\xi_1 \to \sigma\xi_1$ et $\xi_2$, $\sigma\xi_2$, joints
par de petits traits}

\emph{NB} Positions relatives \add{(i.e.\ \uncertain{à conjugaison près})}
d'un $u \in C$ et d'une base $b \subset C$.

(i) $u \in b$. Pour $u$ \add{$= \xi_1$} fixé, il y a 16 telles bases,
correspondant au choix d'un $\xi_1 \in C - E(u)$

(ii) $u \in b'$ (base opposée). Pour $u$ \add{$= \xi_1$} fixé, 16 telles bases,
correspondant au choix d'un $\xi_1 \in C - E(u)$

(iii) $u \notin b \cup b'$. Pour $u$ fixé, il y a 40 telles bases
(\emph{orthogonal} à $u$)
\[
16 + 16 + 40 = 72
\]
\marginal{à gauche de (i)-(iii), reliées à eux par des traits :
$\left.\begin{array}{l} r_b(u) = 1 \\ E(u) \cap b = \emptyset \\
\operatorname{card} E(u) \cap b' = 5 \end{array}\right\rbrace$,
$\left.\begin{array}{l} r_b \cdot u = -1 \\ \operatorname{card} E(u) \cap b =
5 \\ E(u) \cap b' = \emptyset \end{array}\right\rbrace$,
$\left.\begin{array}{l} r_b \cdot u = 0 \\ \operatorname{card} E(u) \cap b = 2
\\ \operatorname{card} E(u) \cap b' = 2 \end{array}\right\rbrace$}
\note{les $\emptyset$ de la marge sont récrits sur d'autres signes ; la
deuxième accolade porte $E(u) \cap b^{*}$ ou $b$ : lecture incertaine. Le
total « 72 » est souligné deux fois}


\end{document}
